% Asistencia de IA (Codex): síntesis del informe y redacción del borrador.
\subsection*{Apartado 6. Tolerancia a fallos (Melanie Muñoz, 2 min)}

El 16 de abril de 2025, Spotify cambió el orden de sus filtros de Envoy en
todas las regiones. El cambio activó un defecto que derribó las instancias
del perímetro. A las 12:20 UTC saltaron las alarmas por una caída del
tráfico. Los reinicios simultáneos y los reintentos de los clientes
dispararon la carga. Como el máximo de memoria de Envoy superaba el límite
de Kubernetes, las nuevas instancias volvían a detenerse. Asia-Pacífico
soportó el fallo con menos carga. Spotify añadió capacidad para romper el
ciclo, corrigió el defecto y el desajuste de memoria, y anunció mejoras en
los despliegues y el monitoreo \cite{spotifyIncident2025}.

\subsection*{Apartado 8. PFC y conclusiones (Melanie Muñoz, 1 min)}

Washington propone separar los archivos de SREI y migrar PostgreSQL a un
servicio administrado. Así, los archivos y la base de datos dejan de depender
del servidor donde corre la aplicación. Los aportes de Jilmar, Elian y
Melanie quedan pendientes de confirmar.

La distribución regional oculta al usuario la ubicación del servicio, pero
un cambio simultáneo puede propagar el mismo fallo a varias regiones.
Spotify separa la búsqueda por solicitud y respuesta del procesamiento
asíncrono de eventos \cite{spotifyIncident2025,spotifyWebApiSearch,khlebnikova2024dataplatform}.
Para analizar CAP o PACELC hay que identificar qué operación y qué estado
compartido se estudian. La caída de Envoy por sí sola describe un fallo del
perímetro, sin establecer cómo se comporta una réplica ante una partición.
